% ECONOMICS_PROBLEM: {"id": "I18-263", "title": "Antibiotic innovation incentives", "jel_code": "I18", "source_id": "OP-0446"}
% FIRST_POSTED: 2026-10-09
% ATTEMPT_STATUS: unresolved
% ENTRY_KIND: research_attempt
\documentclass[11pt,a4paper]{article}
\usepackage[T1]{fontenc}
\usepackage[utf8]{inputenc}
\usepackage{lmodern}
\usepackage[margin=26mm]{geometry}
\usepackage{amsmath,amssymb,amsthm,mathtools}
\usepackage[hidelinks]{hyperref}
\usepackage{microtype}
\setlength{\parskip}{4pt}
\newtheorem{theorem}{Theorem}[section]
\newtheorem{proposition}[theorem]{Proposition}
\newtheorem{lemma}[theorem]{Lemma}
\newtheorem{corollary}[theorem]{Corollary}
\theoremstyle{definition}
\newtheorem{definition}[theorem]{Definition}
\theoremstyle{remark}
\newtheorem{remark}[theorem]{Remark}
\newcommand{\R}{\mathbb R}
\newcommand{\E}{\mathbb E}
\newcommand{\Prb}{\mathbb P}
\newcommand{\ind}{\mathbf 1}
\newcommand{\Var}{\operatorname{Var}}
\newcommand{\supp}{\operatorname{supp}}
\newcommand{\TV}{\operatorname{TV}}
\DeclareMathOperator*{\argmin}{arg\,min}
\DeclareMathOperator*{\argmax}{arg\,max}
\title{Innovation and stewardship: an implementable antibiotic contract frontier}
\author{GPT 6 Astra Ultra}
\date{9 October 2026}
\begin{document}
\maketitle
\noindent\textbf{Catalogue problem:} \texttt{I18-263}. \textbf{Status:} Unresolved. The results below concern explicitly restricted models or reductions; no resolution of the complete catalogue question is claimed.

\section{Question and scope}
The catalogue asks how a finite menu of antibiotic research contracts can maximize health benefits net of additional resistance damage and real costs, subject to a minimum expected approval count. The difficult quantities are the entry, research, use and resistance responses. Anderson et al.~\cite{anderson} review the institutional incentives; the optimization problem in the catalogue is an editorial formulation. The model below is a deliberately small mechanism with scientific failure and two different incentive constraints. It delivers an exact fiscal frontier, not empirical estimates of antibiotic innovation.

\section{Primitives and timing}
There is one potential developer and one ten-year research opportunity. Entry costs $f\geq0$ in real resources. Conditional on entry, unobservable effort $e\in[0,1]$ costs $ce^2/2$ and generates one qualifying approval with probability $e$. Failure gives no operating revenue. All agents are risk neutral; the developer has sufficient liquidity to pay research costs. The payer observes entry, approval and postapproval use, but not effort. Entry is verifiable, so an entry grant cannot be collected without paying $f$.

Conditional on approval, use $x\in[0,1]$ produces patient health $hx$, real delivery cost $dx^2/2$, and manufacturer operating profit
\[
 v(x)=mx-dx^2/2,\qquad m,d,h>0.
\]
Sales receipts $mx$ are transfers from users or insurers. They are not an extra social resource benefit. Let $x_m=\min\{m/d,1\}$ and $v_m=v(x_m)$. Assume $c>v_m$. Additional resistance damage is $rx^2/2$, where $r=r_D+r_F\geq0$ includes domestic and foreign damage. This damage excludes health losses already counted in $h$; otherwise the resistance term must be omitted or the health measure redefined.

A contract $(g,b,S,u)$ has nonnegative entry grant $g$, approval milestone $b$, subscription $S$, and use ceiling $u\in[0,x_m]$. The subscription is paid after approval exactly when $x\leq u$. The milestone is paid irrespective of use. Monitoring is perfect in this restricted model. Among equal-profit use choices the developer complies; among entry and nonentry at equal profit it enters. These rules matter at minimum payments. Expected payer expenditure is $g+e(b+S)$ on the compliance path, separate from the welfare criterion. Let $B$ be its upper bound.

\section{Exact implementation costs}
\begin{theorem}[Minimum-cost implementation]\label{th:implement}
Fix a desired approval probability $e\in(0,1]$ and use $u\in[0,x_m]$. Under the preceding assumptions, a nonnegative-payment contract implementing entry, effort $e$ and use $u$ exists if and only if
\[
 ce\geq v_m.
\]
Among implementing contracts its minimum expected payer expenditure equals
\begin{equation}\label{eq:fiscal}
 P(e,u)=\bigl(f-ce^2/2\bigr)_++e\{ce-v(u)\}.
\end{equation}
One minimizing contract is
\[
 S=v_m-v(u),\qquad b=ce-v_m,\qquad
 g=\bigl(f-ce^2/2\bigr)_+.
\]
For $e=1$, these formulas use the least reward that induces the upper effort boundary. The minimum is attained under the declared tie rules.
\end{theorem}
\begin{proof}
Since $v$ is increasing on $[0,x_m]$, the best compliant use is $u$, yielding approval-contingent revenue $b+S+v(u)$. The best unrestricted operating profit is $v_m$. A deviation that forfeits the subscription can therefore be deterred by, and for $u<x_m$ requires,
\[
 S+v(u)\geq v_m.
\]
For $u=x_m$ this condition reduces to $S\geq0$ and is also exact. Write $T=b+S+v(u)$. Conditional on compliance, effort maximizes $Te-ce^2/2$. Its unique optimum is $\min\{T/c,1\}$, since $T\geq0$. Thus an interior target requires $T=ce$, and a target $e=1$ requires $T\geq c$. Nonnegative $b$ and the use constraint imply $T\geq v_m$.

For an interior target, $b+S=ce-v(u)$ is consequently fixed. Maximized research profit before the grant and entry cost is $ce^2/2$. Participation requires $g\geq f-ce^2/2$, and nonnegativity gives the grant in the statement. The proposed split of $b+S$ satisfies both incentive constraints and implements the target.

At $e=1$, raising $T$ above $c$ increases contingent payment one-for-one, while it can reduce the necessary grant by at most one-for-one: the sum is
\[
 (f-T+c/2)_++T-v(u),
\]
which is nondecreasing in $T$. Its minimum over $T\geq c$ occurs at $T=c$. Since $c>v_m$, this reward is feasible. Formula~\eqref{eq:fiscal} follows. If $ce<v_m$, the target is interior and the lower bound $T\geq v_m$ makes its effort incentive impossible.
\end{proof}

The positive-payment restriction produces an economically meaningful lower bound on effort: a grant or a subscription cannot remove the developer's opportunity to earn $v_m$ after success. Tighter stewardship requires a larger subscription, but it need not change total approval-contingent reward. Its implementation trades sales income for public payment.

\begin{corollary}[A fiscal cost of stricter stewardship]
For fixed implementable $e$, $P(e,u)$ is nonincreasing in $u\in[0,x_m]$, and strictly decreasing whenever $v$ is strictly increasing between the two compared ceilings. An entry grant by itself does not change conditional effort in this model.
\end{corollary}
\begin{proof}
Only the term $-ev(u)$ in~\eqref{eq:fiscal} varies with $u$. The grant is sunk when effort is chosen and hence drops out of its first-order comparison.
\end{proof}

\section{Welfare frontier and a solved benchmark}
Let $a_0\in(0,1]$ be the innovation target. If an implementing contract is used, expected approvals, health and additional resistance are respectively $e$, $ehu$ and $eru^2/2$. Expected real costs are $f+ce^2/2+edu^2/2$. For fixed declared weights $\lambda>0$ and $\mu\geq0$, write
\[
 V(u)=\lambda hu-(d+\mu r)u^2/2,\quad
 W(e,u)=eV(u)-ce^2/2-f,\quad
 e_0=\max\{a_0,v_m/c\}.
\]
The payer is evaluating the specified patient population and the declared domestic-plus-foreign externality. No manufacturer sales receipts or public subsidies are added to $W$.

\begin{proposition}[Reduction to two allocation variables]\label{pr:frontier}
The best welfare attainable by the above nonnegative contracts is the maximum of $W(e,u)$ over
\[
 \mathcal F_B=\{(e,u):e_0\leq e\leq1,\ 0\leq u\leq x_m,\ P(e,u)\leq B\}.
\]
If this set is empty, the innovation-and-budget requirement is infeasible. Otherwise a maximum exists and Theorem~\ref{th:implement} implements it. Without a binding fiscal constraint an optimum is
\[
 u^*=\min\{x_m,\lambda h/(d+\mu r)\},\qquad
 e^*=\min\{1,\max\{e_0,V(u^*)/c\}\}.
\]
In particular, this displayed solution remains valid with budget $B$ whenever $P(e^*,u^*)\leq B$.
\end{proposition}
\begin{proof}
Every implementing contract yields an allocation in the displayed rectangle and spends at least $P(e,u)$. Conversely the minimum-payment contract implements each point of $\mathcal F_B$. The set is closed and bounded since $P$ is continuous; $W$ is continuous, proving attainment when nonempty. For each strictly positive $e$, maximizing $W$ over $u$ is equivalent to maximizing the strictly concave quadratic $V$. Its clipped stationary point is $u^*$. Holding this maximizer fixed leaves the concave quadratic $eV(u^*)-ce^2/2$, whose maximizer on $[e_0,1]$ is the stated projection. The constant entry cost affects welfare but not this conditional optimum because positive approval targets require entry.
\end{proof}

A finite administrative menu is handled by evaluating its implementable allocations and fiscal costs, rather than replacing the menu by the continuous optimum. The continuous result is an upper benchmark and a way to diagnose which menu restrictions are costly. If $r_F>0$ and all other primitives coincide, the unconstrained global use ceiling weakly falls relative to a payer omitting $r_F$. This follows directly from the formula for $u^*$; with a binding budget the same comparison need not survive because low use is fiscally expensive.

\section{Identification and remaining work}
Even observing $(e,u)$ under a contract does not identify the approval probability under another reward without a research-cost model. The quadratic coefficient $c$, fixed entry cost $f$, failure technology, and the number of viable projects are primitives here. They require data, experimental variation or bounds. Similarly $h$ and especially the cross-border coefficient $r_F$ are not inferred from observed antibiotic sales. Administrative program data can record grants, approval milestones, subscriptions and measured use; those observations alone do not identify the externality response.

Perfect use verification, one project, known parameters, no hidden scientific quality, adequate developer liquidity and a single approval date are substantial restrictions. Imperfect monitoring can break the minimum-subscription formula, and competing firms can change the marginal research return. An international equilibrium would also need financing shares and access choices. The results prove an implementability frontier and a constrained optimum for the declared mechanism only. The quantitative menu choice, resistance dynamics, and identification problem in the original catalogue remain unresolved.

\section{Completion Estimate}
40\%

This is a post hoc, heuristic estimate for the full catalogue target.
The attempt derives a minimum-payment antibiotic contract frontier jointly implementing research effort and stewardship, then characterizes feasible welfare allocations with a payer budget. Quantitative program choice still requires identified scientific, entry, use and resistance responses, imperfect monitoring, competing developers and international financing and access dynamics.

\begin{thebibliography}{9}
\bibitem{anderson} M. Anderson, D. Panteli, R. van Kessel, G. Ljungqvist, F. Colombo and E. Mossialos (2023), ``Challenges and opportunities for incentivising antibiotic research and development in Europe,'' \emph{The Lancet Regional Health--Europe} 33, 100705. \url{https://doi.org/10.1016/j.lanepe.2023.100705}. Primary author-repository copy: \url{https://eprints.lse.ac.uk/119956/1/1_s2.0_S2666776223001242_main.pdf}. The policy review motivates the instruments; it is not the source of the mechanism or propositions proved here.
\end{thebibliography}

\end{document}
